\documentclass[10pt]{article}
\usepackage{icmlko}

\icmlmeta{IP 파운데이션 월드모델}{484}{자가 두 개였다 --- 판과 자가 서로 다른 스피어만을 쓰고 있었다}{동률을 배열 순서로 깨는 통계량은 통계량이 아니라 정렬의 함수다. 행을 섞으면 12/12 도메인에서 값이 갈렸다}

\begin{document}

\twocolumn[
\icmlheader

\begin{abstract}
이 실험실의 모든 판정은 한 문장으로 끝난다 --- \emph{``판 $\Delta$ 를 채택 문턱에
댄다.''} 그 문장이 성립하려면 \textbf{판과 문턱이 같은 자로 재져야} 한다.
안 그랬다.

\textbf{판}(\texttt{state/rank\_test.py})은 스피어만을
$\mathrm{argsort}(\mathrm{argsort}(v))$ --- \textbf{서수 순위}로 쟀다. 동률을
\textbf{배열 순서}로 깬다. \textbf{자}(\texttt{runners/ruler890.sp} ·
\texttt{thresh891} 전량)는 처음부터 \textbf{동률 평균}(midrank)이었다. 둘은
동률이 없을 때만 같은 물건이다.

\textbf{동률은 예외가 아니었다.} 챔피언의 예측은 자루 순위 평균이라 같은 값이
흔하다 --- 시장팝업은 채점행 \textbf{126}개에 예측 고유값이 \textbf{15},
영화는 \textbf{406}개에 \textbf{139}뿐이다.

\textbf{그래서 판 $\rho$ 가 자료의 정렬에 의존했다.} 같은 (예측, 라벨) 짝을
스무 번 섞기만 해도 옛 함수는 \textbf{12/12 도메인}에서 값이 갈렸고 폭이 최대
$\mathbf{0.15616}$ 였다. 동률 평균은 \textbf{0/12} · 폭 \textbf{정확히 0} 이다.

\textbf{어느 쪽이 큰지 보기 전에 골랐다.} 사전등록한 규칙은 크기가 아니라 원칙이다
--- \emph{통계량은 행 순서에 의존하면 안 된다.} 그 규칙이 갈랐고, 결과로 정본이
$0.46982 \to \mathbf{0.47034}$ 로 옮겨갔다(12씨앗 짝 $\Delta$
$\mathbf{+0.00051922 \pm 0.00003993}$, 군집 BCa $\lbrack+0.00044,+0.00059\rbrack$,
\textbf{12/12 양수}). \textbf{개선이 아니라 편의 제거다} --- 서수는 동률에서
순위를 무작위로 갈라 상관을 아래로 끈다.

\textbf{움직인 것은 사실상 한 도메인이다.} 시장팝업이 판 기여의 \textbf{98.4\%}
다. 그리고 \textbf{동률이 많다고 움직이는 게 아니다} --- 영화는 동률 비율이
65.8\% 인데 판 기여가 $+1\!\times\!10^{-5}$ 다. 문제는 비율이 아니라 \textbf{동률
블록의 크기}다.

\textbf{그리고 직전 사이클의 진단을 뒤집는다.} 이슈 \#112 는 옛 값과의 차를
\emph{``적합 함수가 GBM 씨앗을 갈아 끼운다''}로 설명했다. \textbf{물리적으로
불가능하다} --- 자루마다 그 값이 자루 색인으로 덮인다(실측: 7 을 넣었는데 자루
32개가 $0,1,\dots,31$). 진짜 원인은 \textbf{이 논문이 고친 그 스피어만}이다.
숫자는 맞았고 \textbf{이름이 틀렸다}: $+0.00052$ 라는 크기는 정확했는데 그것이
잔차에 붙인 이름이었다.

\textbf{자는 안 움직였다} --- 문턱을 짓는 두 성분이 처음부터 동률 평균이었다.
반대로 통일했다면 문턱이 \textbf{+2.8\%} 커졌을 것이다.
\textbf{못 잰 것도 적는다}: 행 군집 자로는 이 $\Delta$ 가 \textbf{판정 불능}이고,
그 자에는 편향이 있다 --- \textbf{재표집이 동률을 스스로 만든다.}
\end{abstract}
]

\section{두 개의 자}

\subsection{배선을 먼저 셌다}

측정 전에 저장소에서 스피어만을 부르는 자리를 전수로 세었다.

\begin{table}[t]\centering\small
\begin{tabular}{lll}
\hline
경로 & 순위 함수 & 자리 \\
\hline
챔피언 판정 & \textbf{서수} & \texttt{rank\_test.py:44} \\
자 (890 · 891) & 동률 평균 & \texttt{ruler890.py:59} \\
837 재적합 & 동률 평균 & \texttt{refit112.py:48} \\
\hline
\end{tabular}
\caption{판과 자가 다른 함수를 쓴다. 세 경로 중 \textbf{판만} 서수다.}
\label{tab:wire}
\end{table}

\claimx{\textbf{판과 자가 서로 다른 스피어만을 쓰고 있었다.} 그래서 이 실험실의
판정 문장 --- \emph{``판 $\Delta$ 를 문턱에 댄다''} --- 은 \textbf{두 개의 자를
섞어 대는 문장}이었다. 문턱을 짓는 \texttt{thresh891} 은 씨앗 성분도 행 성분도
전부 동률 평균으로 계산한다.}
{\texttt{rank\_test.spearman} 을 부르는 경로가 챔피언 채점(\texttt{harness.py:267})
말고 문턱 계산 쪽에도 있으면 이 주장은 약해진다. 전수 grep 으로 세었고
\texttt{thresh891} 의 다섯 통계는 전부 \texttt{ruler890.sp} 를 지난다.}

\subsection{동률은 예외가 아니다}

그림~\ref{fig:one}(a). 채점행 3,775개를 도메인별로 갈라 예측의 고유값 수를
세었다. 시장팝업은 \textbf{126행에 15개}, 영화는 \textbf{406행에 139개}다.
나머지 열 도메인은 동률 비율이 $0.011\sim0.043$ 이다.

\begin{figure*}[t]\centering
\includegraphics[width=\textwidth]{figs/fig1.pdf}
\caption{(a) 도메인별 동률 비율. (b) 12씨앗 전량의 짝 $\Delta$ --- 두 팔이 같은
적합·같은 예측을 공유하므로 완전 짝이다. (c) 판 기여 $\Delta$ 를 $10^{-4}$
단위로. \textbf{값은 전부 산출물 json 에서 읽었다.}}
\label{fig:one}
\end{figure*}

\section{그래서 무엇이 깨졌나}

\subsection{판 $\rho$ 가 자료의 정렬에 의존했다}

사전등록한 첫 규칙이 이것이었다: 측정 러너가 도메인마다 (예측, 라벨) 짝을
\textbf{스무 번 섞어} 두 함수로 $\rho$ 를 재고 \textbf{값의 가짓수}를 센다.

\claimx{\textbf{옛 판 $\rho$ 는 통계량이 아니라 정렬의 함수였다.} 서수 순위는
\textbf{12/12 도메인}에서 값이 갈렸고 폭이 최대 $\mathbf{0.15616}$ 다.
동률 평균은 \textbf{0/12} 이고 폭이 \textbf{정확히 0} 이다(그림~\ref{fig:two}a).}
{한 도메인이라도 동률 평균에서 값이 갈리면 이 주장은 틀린다. 반대로 서수가
어느 도메인에서도 안 갈리면 그 도메인엔 동률이 없다는 뜻이고, 그때는
두 함수가 같은 물건이다.}

\subsection{판정은 크기를 보기 전에 정해 두었다}

사전등록 문서(측정 \textbf{전} 커밋 \texttt{6373373c1} · 21:55:49)에 원칙 셋과
규칙 다섯을 적었다. 규칙 R-4 가 명시한다 --- \emph{``어느 쪽이 큰가는 판정에
안 쓴다. 정본이 얼마나 옮겨지는지는 결과로만 보고한다.''}
팔 B 를 못 고르는 실패 조건도 하나만 미리 적었다: \textbf{887형 중립화}
(midrank 가 NaN 하나로 열을 통째 NaN 으로 만드는 경로). NaN 을 \textbf{심어서}
확인했고 --- 마스크 없이는 실제로 \texttt{nan} 이 나온다 ---
\texttt{harness.py:265} 의 \texttt{ok} 마스크가 막는다. 발화하지 않았다.

\section{그 결정이 정본을 얼마나 옮기나}

두 팔은 \textbf{같은 적합·같은 예측 벡터}를 공유한다. 채점 함수만 다르므로
완전 짝이고, 씨앗마다 한 번만 적합했다.

\begin{table}[t]\centering\small
\begin{tabular}{lr}
\hline
팔 A (서수 · 옛 정본) & $0.46982330 \pm 0.0020371$ \\
팔 B (동률 평균 · \textbf{새 정본}) & $\mathbf{0.47034252} \pm 0.0020895$ \\
\hline
짝 $\Delta$ (12씨앗) & $+0.00051922 \pm 0.00003993$ \\
$\Delta>0$ 인 씨앗 & \textbf{12 / 12} \\
씨앗 짝 BCa 95\% & $\lbrack+0.00044404,\ +0.00059359\rbrack$ \\
행 군집 BCa 95\% & $\lbrack-0.00290,\ +0.00342\rbrack$ \\
\hline
\end{tabular}
\caption{12도메인 · 유보 3,775 · 씨앗 0--11. 팔 A 는 챔피언 경로의 기록값을
$5.6\times10^{-17}$ 안에서 재현한다.}
\label{tab:board}
\end{table}

\claimx{\textbf{정본이 $0.46982 \to 0.47034$ 로 옮겨간다.} 짝 $\Delta$
$+0.00051922$ · \textbf{12/12 양수} · 씨앗 짝 BCa 가 0 을 안 문다.
새 SE 의 $0.86$ 배 · 채택 문턱의 $0.147$ 배다.
\textbf{개선 주장이 아니다} --- 서수는 동률에서 순위를 무작위로 갈라 상관을
\emph{아래로} 끄는 편의가 있었고, 그 편의를 뺀 것이다. 손잡이가 아니라
\textbf{통계량의 정의}가 바뀌었으므로 ``첫 양수를 채택하지 않는다''는 조항의
대상이 아니다.}
{도메인 기여의 합이 짝 $\Delta$ 와 안 맞으면(실측 차 $1.8\times10^{-17}$),
또는 팔 A 가 챔피언 기록값을 $10^{-12}$ 안에서 재현하지 못하면 배선이 다른 것이다.}

\subsection{못 잰 것을 못 쟀다고 적는다}

행 군집 BCa 는 \textbf{판정 불능}이다($\lbrack-0.0029,+0.0034\rbrack$). 그리고
\textbf{그 자에는 편향이 있다} --- 행을 재표집하면 \textbf{동률이 새로 생기고},
이 논문이 재는 대비가 바로 동률 처리이므로 폭이 위쪽으로 부푼다. 노트 890 이
이미 자기 코드 주석에 그 사실을 등록해 두었다.

\claimx{\textbf{``행이 달랐다면 두 정의가 갈렸을까''는 아직 모른다.} 씨앗 짝
자는 승이라고 말하고 행 군집 자는 판정 불능이라고 말한다. \textbf{둘 다 적는다.}
결정은 어느 쪽 자에도 안 걸려 있다 --- 원칙으로 골랐기 때문이다.}
{재표집이 만든 동률과 자료가 원래 가진 동률을 가르는 자를 만들면 이 칸이
채워진다. 이 사이클에서는 안 만들었다.}

\subsection{움직인 것은 사실상 한 도메인이다}

그림~\ref{fig:one}(c). 시장팝업 하나가 판 기여의 \textbf{98.4\%} 다.
아이돌이 10.2\%, 펀딩이 $-18.7\%$(유일하게 반대로 간다), 나머지 아홉은 합쳐서
1\% 미만이다.

\claimx{\textbf{동률이 많다고 움직이는 게 아니다.} 영화는 동률 비율이
\textbf{65.8\%} 인데 판 기여가 $+1.1\times10^{-5}$ 로 사실상 0 이다.
시장팝업(88.1\%)과 갈리는 것은 비율이 아니라 \textbf{동률 블록의 크기}다 ---
$126/15 = 8.4$ 대 $406/139 = 2.9$. 무작위 타이브레이크가 상관을 끄는 정도는
블록 안에서 섞이는 순위의 수에 달렸다.}
{블록 크기가 비슷한데 기여가 크게 다른 도메인 짝이 나오면 이 설명은 부족하다.
현재 열두 도메인으로는 블록 크기가 큰 사례가 둘뿐이라 \textbf{표본이 둘이다}.}

\section{잔차의 크기는 맞았고 이름이 틀렸다}

옛 정본과 그 이전 값의 차를 12씨앗 전량으로 갈랐다. 네 조합을 \textbf{따로
적합해} 갈랐다 --- 직전 사이클은 도메인을 갈라 잔차에 이름을 붙였고, 그
분해는 구성상 항등식이라 검증이 아니었다(그 러너가 자기 파일에 그렇게 적어
두었는데 인쇄물은 그 단서를 떼고 단정했다).

\[
\underbrace{+0.00117370}_{\text{총}}
=\underbrace{+0.00065447}_{\text{채점 배치}}
+\underbrace{+0.00051922}_{\text{스피어만 구현}}
\]

합 검산은 \textbf{정확히 0} 이고 두 성분 모두 \textbf{12/12} 씨앗에서 양수다.
첫 항은 노트 855 가 적어 둔 $+0.00065$ 와 소수 다섯째 자리까지 같다.
\textbf{둘째 항의 크기도 직전 사이클이 적은 값과 같다.} 틀린 것은 그
크기가 아니라 \textbf{그것에 붙인 인과}였다.

\claimx{\textbf{`적합 함수가 GBM 씨앗을 갈아 끼운다'는 기전은 물리적으로
불가능하다.} 자루마다 그 값이 \textbf{자루 색인}으로 덮인다. 일부러 어긋나게
주고(7) 자루를 찍으니 $0,1,\dots,31$ 이었다. 그 인과가 여덟 자리에 인쇄돼
배포됐고, 진짜 원인은 \textbf{오늘의 정본 자신의 결함}이었다.}
{자루의 \texttt{random\_state} 가 주입한 값과 같게 나오면 이 주장은 틀린다.
단, 같은 함수가 설정하는 \emph{다른} 속성(자루 부트스트랩 rng)은 무연산이
\textbf{아니고}, 그것은 양쪽 경로가 똑같이 한다.}

\paragraph{정정 표시.} 이 절이 인용하는 $0.4710$ 은 \textbf{837 시대 값}이고
판정에서 내려온 수다. 837 규약(라벨 배치 채점 + 동률 평균)으로는 \emph{오늘도}
재현되지만(0.4709970), 챔피언 규약으로는 안 나온다. 이 논문에서 그 수는
\textbf{분해의 대상}으로만 쓰이고 어떤 판정에도 안 들어간다.

\section{자는 움직였나}

\begin{figure}[t]\centering
\includegraphics[width=\columnwidth]{figs/fig2.pdf}
\caption{(a) 행을 스무 번 섞었을 때 나온 $\rho$ 의 폭(로그 축). 동률 평균은
정확히 0 이라 그려지지 않는다. (b) 두 순위 함수 아래의 채택 문턱.}
\label{fig:two}
\end{figure}

사전등록에 \textbf{반증 가능한 예측}으로 적었다 --- \emph{``문턱을 짓는 두 성분이
전부 동률 평균이므로 팔 B 는 문턱을 안 움직인다.''} 문턱 러너의 적합 캐시를
\textbf{읽기만} 해서(재적합 0회) 같은 난수·같은 뽑기 순서로 두 함수를 돌렸다.

\begin{table}[t]\centering\small
\begin{tabular}{lrrr}
\hline
 & 씨앗 & 행 짝 & 문턱 \\
\hline
동률 평균 & 0.00099 & 0.00146 & \textbf{0.0035242} \\
서수 & 0.00105 & 0.00148 & 0.0036223 \\
\hline
\end{tabular}
\caption{성분은 다섯째 자리까지 기존 값을 정확히 재현한다.}
\label{tab:thr}
\end{table}

\claimx{\textbf{팔 B 는 자를 안 움직인다} --- 자는 처음부터 동률 평균이었다.
반대로 \textbf{서수로 통일했다면 문턱이 $+2.8\%$ 커진다}(0.0035242 $\to$
0.0036223). 즉 이 선택은 판만 옮기는 것이 아니라 \emph{어느 쪽으로 통일하느냐에
따라 문턱도 움직였을} 선택이었다.}
{문턱 러너의 다섯 통계 중 하나라도 서수 함수를 지나면 이 주장은 틀린다.
캐시를 읽어 두 함수로 각각 다시 계산했고 동률 평균 쪽이 기존 성분을 재현했다.}

\paragraph{곁 발견 --- 문턱은 반올림한 수로 지어져 있었다.} 인쇄된 문턱
$0.00353$ 은 \textbf{다섯 자리로 반올림한 성분}을 $2\sqrt{a^2+b^2}$ 에 넣어 나온
값이다($0.0035280$). 반올림 없이는 $0.0035242$ 다. 차 $4.0\times10^{-6}$ 로
어떤 판정도 안 바뀌고 방어 대역 안이지만, \textbf{판정을 가르는 상수가 반올림한
입력 위에 서 있었다}는 사실 자체가 기록될 값어치가 있다.

\section{한계}

\begin{itemize}\itemsep2pt
\item 행 군집 자는 \textbf{판정 불능}이고 편향 방향이 알려져 있다. 편향 없는
행 자를 안 만들었다.
\item 같은 함수를 쓰는 \textbf{다른 열여덟 모듈}의 산출물을 다시 안 돌렸다.
그 값들은 전부 서수 시대의 수다.
\item 옛 함수는 비유한 값을 맨 뒤로 밀어 \textbf{조용히 숫자를 냈다}. 새 함수는
막는다 --- 그 가드가 열여덟 모듈 중 어디서 발화하는지 \textbf{안 세었다}.
발화하는 곳이 있으면 그건 원래 틀린 값이었다는 뜻이다.
\item 블록 크기 설명의 \textbf{표본이 둘}이다(시장팝업 · 영화).
\item 정본을 나르는 자리 일곱 곳은 다른 작업자가 동시에 만지는 파일이라
\textbf{내가 못 고쳤다}. 정본 대조기가 설계대로 소리를 냈고 이슈로 열었다.
\end{itemize}

\end{document}
